\documentclass[11pt]{article}
% Uses the canonical latex-docs semantic style in-repository. The local
% lightweight fallback is for independent compilation of this distribution.
\IfFileExists{technical-user-manual.sty}{\usepackage{technical-user-manual}}{\usepackage{../support/aaoa-standalone}}
\usepackage{mathtools}
\setDocTitle{Chapter 6: Trees}
\setDocSubtitle{Guided reading, derivations and original exercises}
\setDocVersion{1.0}
\setDocStatus{Educational reference / unapproved working material}
\setDocOwner{Algorithm Engineering}
\setDocAuthor{Original companion material}
\setDocSummary{Coverage of primary book Chapter 6 with independently authored examples, proof obligations and lab transfer.}
\providecommand{\Expect}{\mathbb{E}}
\providecommand{\Prob}{\mathbb{P}}
\providecommand{\Var}{\operatorname{Var}}
\providecommand{\bigoh}{\mathcal{O}}
\begin{document}
\makedoctocpage
\section{Reading map and prerequisite check}
\noindent\textbf{Primary source sections:} 6.1 Binary Trees; 6.2 Forests and Trees; 6.3 Equivalences; 6.4 Properties; 6.5 Tree Algorithms; 6.6 BSTs; 6.7 Average Path Length in Catalan Trees; 6.8 Path Length in BSTs; 6.9 Additive Parameters; 6.10 Height; 6.11 Summary; 6.12 Lagrange Inversion; 6.13 Rooted Unordered Trees; 6.14 Labelled Trees; 6.15 Other Types (printed pp. 258--331).\par
\noindent\textbf{Prerequisites:} Completed preceding mathematical-method chapters or equivalent demonstrated mastery.\par
\section{Learning outcomes}
\begin{enumerate}
\item Distinguish uniform Catalan shape from random BST insertion law.
\item Derive an additive-tree-parameter recurrence.
\item Measure internal path length and height with consistent definitions.
\end{enumerate}
\section{Guided technical analysis}
\subsection{What is a random tree?}
A uniformly chosen tree shape among all Catalan shapes of size $n$ is not the same distribution as a BST from a uniformly random key insertion order. For the latter, each rank has probability $1/n$ of being root; recursive subtrees inherit the random-permutation model.
\subsection{Expected BST internal path length}
With root depth zero and $P_0=P_1=0$, the expected total depth $P_n$ satisfies
\[
P_n=n-1+\frac{2}{n}\sum_{k=0}^{n-1}P_k,
\qquad P_n=2(n+1)H_n-4n.
\]
This identity describes \emph{total} internal path length, not average node depth. Divide by $n$ only when the desired statistic is mean depth per node.
\subsection{Height is different}
Height is a maximum, so additive-parameter recurrences do not apply directly. Define whether empty-tree height is $-1$ or zero and keep that convention in all tests.
\section{Independent problem set}
\begin{enumerate}
\item Enumerate all insertion permutations for $n=4$, compute tree path length and verify the mean.
\item Compare simulated Catalan-uniform and random-BST height histograms at the same size.
\item Prove that internal path length equals the sum over nodes of their depths by two independent traversals.
\end{enumerate}
\section{Laboratory transfer}
\textbf{Required laboratory:} LAB-PY03: exact tree shape enumeration and BST Monte Carlo.\par
\noindent\textbf{Minimum recorded evidence:} repository commit or source hash; test inputs/seeds; actual operation counts; predicted values; one explicit model limitation.\par
\section{Mastery gate and self-review}
A report labels the tree distribution, node counting, depth convention and sample size.\par
\begin{enumerate}\item Can I define the random object/input exactly? \item Is my identity exact, asymptotic, or empirical? \item Which base cases and counterexamples did I check? \item What would falsify my performance hypothesis?\end{enumerate}
\section{Suggested notebook structure}
Use five named sections: model and invariants; derivation; deterministic small instances; randomized replication; discrepancy and revised hypothesis. Keep a proof separate from measurement results.

\end{document}
